\chapter{Discussion}
\label{ch:discussion}

This chapter turns the design space and the measurements into a
recommended configuration, states which theoretical predictions the design
could test, places the findings against the literature, and sets out the
threats to validity.

\section{From Design Space to Recommendation}
\label{sec:synthesis}

\subsection{The Structure of the Problem}

The FBA design space contains almost no free improvements: each block
offers options that shift competition from one variable to another rather
than removing it. The one genuine structural gain is the mechanism itself,
discrete accumulation plus uniform pricing, which removes the value of a
marginal speed advantage without requiring anything from participants;
everything else is a choice about which residual distortion to accept. The
blocks also interact and cannot be optimised independently: an open book
raises the value of a randomised close, whereas a sealed-bid design makes a
fixed boundary defensible, and a short interval in a thin instrument
produces empty batches and degenerates into a delayed continuous market.

\subsection{A Recommended Baseline}

\Cref{tab:recommendation} states a baseline for a liquid instrument on a
centralised venue. The recommendation follows from the theory alone; the
final column records which measurement bears on each choice and states
plainly where the corresponding comparison was not run.

\begin{table}[htbp]
\centering
\caption[Recommended baseline configuration]{A recommended baseline FBA
configuration for a liquid instrument on a centralised venue.}
\label{tab:recommendation}
\small
\begin{tabularx}{\textwidth}{@{}p{2.3cm} p{3.0cm} L p{2.4cm}@{}}
\toprule
\textbf{Block} & \textbf{Recommendation} & \textbf{Reasoning} &
\textbf{Evidence}\\
\midrule
Interval & 50--250\,ms, with a short randomised extension &
Long enough that the ratio of latency dispersion to $\tau$ is
negligible; short enough that delay is immaterial to any participant
other than a market maker managing inventory & Interval sweep not run
(\Cref{tab:sweep}); the one interval measured, $\tau=1$\,s, is already
well above this range and \FbaEmptyBatchShare\% empty even there\\
Disclosure & Indicative clearing price and aggregate imbalance; no
order-level book & Retains the coordination benefit of transparency
without exposing individual orders; the regime equity auctions already
use at scale & Not testable here (P7)\\
Segregation & Unified pool by default; intent segregation only where
maker flow is demonstrably being adversely selected & Segregation splits
liquidity, which is a certain cost against an uncertain benefit; the
DFBA argument is strongest where the unified pool shows persistent
maker-versus-maker crossing & \Cref{tab:res21} price impact\\
Rationing & Pro rata with a small random component & Removes the timing
channel, which is the one that would undo the mechanism; the random
component blunts the size-inflation incentive without imposing full
execution uncertainty & Rationing-rule comparison not run
(\Cref{tab:rationingres}); the rule actually measured is
price-time-within-the-batch, not this recommendation, so the two should
not be conflated\\
Settlement & Net settlement with per-fill audit records retained;
roll-over residual with standard time-in-force & Netting is a cost
saving with no incentive consequence provided the audit trail is kept
separately & Residual share, \Cref{tab:res23}\\
\bottomrule
\end{tabularx}
\end{table}

\subsection{How the Recommendation Changes by Venue}

\textbf{Thin instruments.} The binding constraint is batch thickness, not
speed neutralisation: a longer interval, or a volume-triggered close, keeps
batches populated at the cost of longer delay. The degeneracy diagnostic is
the empty-batch share, measured only at $\tau=1$\,s since the sweep tracing
it as a function of $\tau$ was not run.

\textbf{Decentralised venues.} The interval is not a free parameter,
because settlement is already discrete and the natural boundary is the
block, and disclosure is constrained since submitted intents are typically
public. In exchange the venue gains a multi-asset clearing formulation with
coherent cross-rates across all pairs in a batch, removing cross-pair
arbitrage as a rent rather than paying arbitrageurs to enforce consistency
\citep{walther2021multitoken,ramsay2023speedex}.

\textbf{Venues with a large retail segment.} The Glosten--Milgrom argument
for segregation is strongest here, because the uninformed subset is both
large and identifiable. The countervailing consideration is regulatory:
participants in different segments receive different prices for the same
asset at the same instant, which is difficult to reconcile with a uniform
best-execution standard.

\subsection{The Unresolved Question: Maker and Taker}

The maker--taker distinction has no natural definition under simultaneous
uniform-price clearing, and fee and rebate schedules are a substantial
share of exchange revenue, so a venue that adopts batching without
resolving the definition will find its rebate schedule subsidising
whichever behaviour its chosen proxy tracks. Of the four candidate
redefinitions, persistence across batches has the strongest economic claim:
it alone measures standing liquidity provision rather than an instantaneous
order property, and an order built solely to capture the rebate cannot
satisfy it. The cost is that it requires tracking order lineage across
batches, which only works if unmatched remainders roll over rather than
being cancelled.

\section{What This Design Can and Cannot Establish}
\label{sec:testable2}

The nine predictions carried through this thesis split by testability, and
the split bounds every claim made.

\textbf{Mechanical predictions are established.} P1 (zero intra-interval
dispersion), P4 (the delay lower bound), P8 (the effect of $\tau$ on batch
composition), and P9 (positive surplus relative to submitted limits) follow
from the clearing rule applied to a given order set, and where the
simulation confirms them by construction rather than by measurement, the
results say so.

\textbf{One prediction is partly reachable.} P2, that batching reduces the
adverse selection component of the effective spread, is testable in its
mechanical part --- given this flow, does clearing it in batches produce
less price impact than clearing it continuously? --- but not its
behavioural part, since liquidity providers in the replay cannot requote in
response to knowing they will not be sniped. A measured reduction in price
impact is therefore a lower bound on the effect theory predicts.

\textbf{Behavioural predictions are out of reach.} P3, P5, P6 and P7 all
require participants who can adapt: quoted spreads narrow because market
makers requote, size inflation occurs because participants respond to pro
rata, boundary concentration because they respond to a predictable gate,
bid shading because they respond to price uncertainty. None of these
responses can occur in a replay of orders submitted under a different
mechanism; the simulation provides the mechanical baseline against which
they would have to be measured. A replay is the right instrument for
isolating what a mechanism does to a given order flow, and the wrong one
for what it does to an equilibrium.

\section{Positioning Relative to the Literature}
\label{sec:positioning}

Against \citet{budish2015hft}, this thesis takes the mechanism as given and
asks how it is built: their contribution is the equilibrium argument, the
contribution here the enumeration of what remains undetermined once it is
accepted, and the observation that several remaining choices can undo part
of the benefit. Price-time rationing inside the batch is the clearest
example: a venue can implement frequent batch auctions exactly as specified
and still have a speed race.

Against \citet{aquilina2022quantifying}, the relationship is complementary:
they measure what the continuous mechanism costs using field data, this
thesis what an alternative mechanism does to the same flow, and their
estimate is the benchmark against which any measured reduction in adverse
selection should be scaled, allowing for their equity market's different
participant population and latency distribution.

Against \citet{lee2026taiwan}, the relationship is methodological: their
real transition gives high external validity but limited internal
comparability, since everything moved at once, while this thesis has perfect
internal comparability, the flow being identical by construction, but
limited external validity, since the flow cannot respond. Together they bound
the answer from two sides.

Against the decentralised exchange literature
\citep{walther2021multitoken,ramsay2023speedex,dutta2025cow}, which solves
the multi-asset clearing problem in practice but evaluates it mostly in
terms of throughput, gas cost and solver incentives, the contribution is to
import the clearing formulation into a microstructure framing, making
visible that those systems are instances of one mechanism rather than a
separate technology.

\section{Threats to Validity}
\label{sec:validity}

\paragraph{Internal validity.} The main threat is that a measured
difference reflects the implementation rather than the mechanism. The two
engines share only an order and trade representation, not a matching
interface, so comparability rests on identical input and one shared metrics
implementation, reinforced by the test layers and by reporting throughput
and latency as explicit controls. The $\tau \to 0$ convergence check would
be the strongest single test, having a known answer and exercising the
entire pipeline, but it was not run.

\paragraph{Construct validity.} Several metrics do not transfer unchanged
from a continuous to a batch setting. The quoted spread has no direct FBA
analogue and is proxied by the marginal unexecuted orders of the batch.
More consequentially, the FBA has no taker/maker distinction, so its
effective and realized spread need a substitute direction before they can
be signed; a tick rule sourced from the CDA's own price series is used
because one from the FBA's own would be circular. This is a proxy for the
missing aggressor flag, not a recovery of it --- tick-rule classification is
imperfect even with true trade-level data, and borrowing the other engine's
tick is a further remove --- so the near-zero FBA spread is evidence against
a one-sided cost under this construction, not a precise cost estimate.

\paragraph{External validity.} The data come from one venue, and a crypto
perpetuals exchange has a participant population, fee structure and latency
profile that differ from an equity market, so magnitudes should not be
transported. The sample is one instrument over one continuous month, so
there is no per-session or per-instrument breakdown, only day-by-day
stability within the month. And the flow is fixed: the mechanism cannot
change behaviour, so any equilibrium effect is absent by construction.

\paragraph{Data validity.} The dataset is captured at a single
non-validating node \citep{albers2026openbook}, and any message the node
missed or saw in a different order than the venue's own engine propagates
into both arms of the comparison. Because it does so identically it is
largely differenced out of the paired comparison, but it affects absolute
levels, and a CDA fidelity check against the venue's own trades would be
informative rather than redundant if it existed.
